\section{Conclusion}
\label{sec:conclusion}

We examined whether large language models exploit an ``answer anyway'' loophole when they are uncertain or when a request is impossible to satisfy exactly. Across \simpleqa, \truthfulqa, and impossible \halogen false-presupposition prompts, the strongest effect appears on impossible structured requests. When abstention is forbidden, both tested models move to universal attempted compliance, and most of those attempts are fabricated.

The broader claim requires nuance. On factual QA, loophole behavior is strongly model-dependent. \gptfivemini shows a clear abstention-versus-bluffing tradeoff, while \gptfourone mostly answers regardless of instruction. The main takeaway is that surface compliance can displace epistemic honesty, especially when prompts remove the abstention option.

Future work should test more model families, include broader abstention-focused benchmarks, and compare prompt-only mitigation against external answer validation. That combination would help determine whether the loophole is best addressed by training, interface design, or system-level checks.
